\mkchapter{Evaluation}
\label{ch:evaluation}

The preceding chapters built an encoding of B proof obligations into SMT-LIB and established its correctness.
This chapter evaluates its \emph{effectiveness} by asking whether the higher-order, representation-aware design of \Cref{ch:loosening,ch:beer} helps an SMT solver discharge industrial proof obligations relative to the first-order encodings used by current tools.
In one controlled campaign, we encoded every goal in a large public corpus of Atelier~B proof obligations~\cite{deharbe2022pog} with both \beer{} and Atelier~B's ppTrans translator, solved the resulting scripts with cvc5 under identical settings, and compared the outcomes goal by goal: \bmGoals{} goals from \bmFiles{} files.

Its outcome has two parts.
As an \emph{encoding}, the higher-order design stands within a few points of the industrial baseline on the goals that both pipelines answer, leads on arithmetic and cardinality, matches the baseline on higher-order goals, and proves tens of thousands of goals that the baseline cannot.
As an end-to-end \emph{tool}, however, \beer{} remains well behind because its encoder fails to cover or finish many files.
Outperforming existing tools is not what \beer{} was built for: it exists to exercise the higher-order capabilities of modern SMT solvers on industrial proof obligations, and the campaign measures where those capabilities stand today.
The chapter substantiates both parts and keeps them distinct: \Cref{sec:evaluation-goals} fixes the questions and the artifacts under test, \Cref{sec:benchmark-corpus,sec:experimental-setup} describe the corpus, the protocol, and the two scoring rules that separate them, \Cref{sec:beer-results} presents the results, and \Cref{sec:threats-to-validity} examines the threats to their validity.

\section{Evaluation goals}
\label{sec:evaluation-goals}

The object of this chapter is the automated pipeline as a whole: a B proof obligation goes in, an SMT script comes out, and a solver either discharges the goal or does not.
Its design rests on two decisions defended in the preceding chapters---representing sets by characteristic predicates and functions by function values rather than by a first-order membership axiomatization (\Cref{ch:beer}), and reconciling the representations this creates through type loosening (\Cref{ch:loosening}).
Both decisions were justified there on semantic and proof-theoretic grounds; whether they help an actual solver is an empirical question, and this chapter answers it through four concrete assessments:

\begin{description}[font=\normalfont\emph]
  \item[Effectiveness:] of all proof obligations in the corpus, how many does the pipeline discharge end to end, and what stops it on the rest?
  \item[Quality:] when both encoders produce a script and the solver reaches a verdict on both, which encoding serves the solver better?
  \item[Structure:] which classes of goals favor which encoding---and do the design choices of \Cref{ch:loosening,ch:beer} pay off where they were predicted to?
  \item[Cost:] what do the higher-order scripts cost in solver time and script size?
\end{description}

\paragraph*{Baseline.}
The comparison point is \emph{ppTrans}, the translator distributed with Atelier~B~\cite{lecomte2024proving} and the industrial reference for discharging B proof obligations with SMT solvers~\cite{deharbe2014integrating}.
It occupies the opposite pole of the design space surveyed in \Cref{ch:introduction}: set-theoretic constructs are compiled away into first-order formulas over simple membership predicates, so no set, function, or relation survives as a first-class value.
Running both translators on the same goals, with the same solver under identical settings, therefore compares the two encoding philosophies rather than two solvers.

\paragraph*{Measured artifact.}
The questions above instead concern the encoding \emph{design} at full scale, so we measure the variant that implements the encoding rules of \Cref{ch:beer}, extended to the constructs the verified fragment excludes and optimized more aggressively.
The encoder under test is therefore \textsc{Zero-Proof} \beer{} (\Cref{sec:beer-zero-proof}).
The distinction is kept explicit throughout: the measured binary shares the design, but its implementation is not fully covered by the soundness theorem of \Cref{sec:beer-soundness}.

\paragraph*{Solver.}
Both scripts are discharged by cvc5~\cite{barbosa2022cvc5}, the reference implementation of higher-order SMT solving~\cite{barbosa2019hosmt}.

\begin{note}
The interactive route through \barrel{} is not evaluated here: a proof-assistant integration is exercised by users rather than by a corpus sweep, and \Cref{ch:barrel} discusses it qualitatively.
\end{note}
\section{Benchmark corpus}
\label{sec:benchmark-corpus}

\section{Experimental setup}
\label{sec:experimental-setup}

\paragraph*{Protocol.}
Every \texttt{.pog} file is translated twice, by \beer{} (revision \texttt{83fda63}, built with Lean~4.27.0) and by the ppTrans translator of Atelier~B~24.04.2.
Both emit exactly one \smtinl{check-sat} per goal, in the order the goals appear in the file, and each goal becomes a self-contained script solved in its own process---no incremental mode, no state shared between goals---so answer \(k\) on one side and answer \(k\) on the other describe the same obligation.
Files where the two goal counts differ are excluded from the join rather than silently misaligned.
Each script is discharged by cvc5~1.3.4 with model-based quantifier instantiation (\texttt{-{}-mbqi}; \Cref{sec:ar-smt}) and a per-\smtinl{check-sat} budget of \(5\)\,s (\texttt{-{}-tlimit-per=5000}), the configuration under which \beer{}'s scripts are intended to run (\Cref{sec:beer-po-running-example}).
The budget is soft---cvc5 checks it between phases---so each solver process is additionally killed at a \(300\)\,s wall; the encoders themselves run under a \(1\,800\)\,s budget per file.
\beer{}'s shipped prelude requests unsat cores, which ppTrans never does and which costs cvc5 real time, so the option is stripped and both encodings are solved under identical solver settings.

\paragraph*{Hardware.}
The campaign ran on the \texttt{gros} cluster of Grid'5000\footnote{\url{https://www.grid5000.fr}} (Nancy): identical nodes with one Intel Xeon Gold 5220 processor, running one solver instance per physical core.
All timings quoted below are from this configuration.

\paragraph*{Scored population.}
Three files had their sweep stopped partway and contribute only the goals actually measured; in total, \bmFiles{} of the corpus's 5\,421 files contribute at least one scored goal, for \bmGoals{} goals.
Deliberately \emph{not} set aside: the goals \beer{}'s encoder produced no script for, and the goals on which a solver was killed at the wall.
The first is encoder coverage---dominated by the files using Atelier~B's \texttt{Struct} records, a construct specific to Atelier~B rather than to the B language of the B-Book~\cite{abrial1996bbook}, which \beer{} does not aim to support---and the second is a failure to discharge within the budget; both are real outcomes of running the tool.

\paragraph*{Two scorings.}
Every result in this chapter uses one of two scoring rules, and confusing them changes every conclusion, so they are fixed here once.

\begin{description}[font=\normalfont\itshape]
  \item[strict] All \bmGoals{} goals; a goal counts as \emph{proved} iff cvc5 answered \smtinl{unsat}.
    A goal \beer{} could not encode, and a goal ppTrans burned the \(300\)\,s wall on, both count as failures to prove.
    This is what a user of either tool experiences end to end, and it is the headline scoring.
  \item[verdict-only] The \bmHead{} goals for which \emph{both} pipelines obtained a solver verdict (\smtinl{unsat} or \smtinl{unknown}).
    This isolates the quality of an encoding from the robustness and coverage of the encoder that produced it, at the cost of discarding \bmDropped{} goals---\(37\,\%\) of the corpus---and discarding them asymmetrically, since each side loses a different population.
\end{description}

\Cref{sec:beer-results-global} reports both scorings, because the gap between them is itself a result.
The structural analyses of \Cref{sec:beer-results-where,sec:beer-results-alone} are verdict-only throughout: they ask which encoding suits which kind of goal, and a goal that never reached the solver cannot answer that.
Every figure names its scoring, and the reference tables closing the chapter carry both side by side.

\section{BEer results}
\label{sec:beer-results}

\section{Threats to validity}
\label{sec:threats-to-validity}

The campaign controls the solver, its options, the hardware, and the per-goal budget, but its conclusions remain conditional on those choices.
Only cvc5~1.3.4 was evaluated, with model-based quantifier instantiation and a \(5\)\,s internal budget.
Different solver versions, higher-order solvers, quantifier strategies, or budgets could change both the absolute success rates and the relative ranking of the encodings.
The timing results in particular describe this configuration rather than an intrinsic cost of either encoding.

The corpus is large and industrial, but it is a single anonymized collection generated by one toolchain.
Its projects contribute very different numbers of goals, and its near-duplicate components carry disproportionate weight; the file-level discussion in \Cref{sec:beer-results-alone} identifies the clearest instance.
The observed subject and construct effects may therefore not generalize to other B developments or proof-obligation generators.
Moreover, the structural classes are descriptive: subject classes are mutually exclusive, construct classes overlap, and neither classification by itself establishes that a construct causes a solver outcome.

The two scoring rules address different selection effects but cannot eliminate them.
Strict scoring measures the end-to-end tools while combining encoder failures with solver failures.
Verdict-only scoring better isolates the emitted encodings, but removes \bmDropped{} goals asymmetrically and therefore does not provide an unbiased estimate of how the encodings would compare on the full corpus.
For this reason, the chapter reports both and avoids transferring conclusions from one population to the other.

Finally, the measured \textsc{Zero-Proof} \beer{} binary implements the verified design but extends and optimizes it beyond the fragment covered by the soundness theorem.
The absence of a \smtinl{sat}/\smtinl{unsat} disagreement with ppTrans is reassuring, but it is not a correctness argument: both translators and the goal-alignment machinery used in the campaign are unverified code.

% Reference tables closing the Evaluation chapter: every breakdown of
% \Cref{sec:beer-results} under both scorings, so that no figure's restriction
% hides the other reading.  Deliberately without a heading or introduction ---
% they are lookup material, and the chapter's prose quotes what matters.
\FloatBarrier

\begin{table}[t]
\centering\small
\setlength{\tabcolsep}{3.3pt}
\pgfplotstabletypeset[
  columns={label,n,beer,pptrans,delta,noscript,pptkilled,nh,beerh,pptransh},
  columns/label/.style={string type, column name={class}, column type={l}},
  columns/nh/.style={column name={goals}, int detect,
                     /pgf/number format/1000 sep={\,}},
  columns/n/.style={column name={goals}, int detect,
                    /pgf/number format/1000 sep={\,}},
  columns/beer/.style={column name={\beer{} \%}, fixed, fixed zerofill,
                       precision=1},
  columns/pptrans/.style={column name={ppTrans \%}, fixed, fixed zerofill,
                          precision=1},
  columns/delta/.style={column name={$\Delta$}, fixed, fixed zerofill,
                        precision=1},
  columns/noscript/.style={column name={no script \%}, fixed, fixed zerofill,
                           precision=1},
  columns/pptkilled/.style={column name={ppT killed \%}, fixed, fixed zerofill,
                            precision=1},
  columns/beerh/.style={column name={\beer{} \%}, fixed, fixed zerofill,
                        precision=1},
  columns/pptransh/.style={column name={ppTrans \%}, fixed, fixed zerofill,
                           precision=1},
  every head row/.style={
    before row={\toprule & \multicolumn{6}{c}{strict}
                & \multicolumn{3}{c}{verdict-only} \\
                \cmidrule(lr){2-7}\cmidrule(l){8-10}},
    after row=\midrule},
  every last row/.style={after row=\bottomrule},
]{\bysubj}
\caption{Strict and verdict-only success rates by subject class, with the two
  failure modes that separate them.}
\label{tab:eval-subject}
\end{table}

\begin{table}[t]
\centering\small
\setlength{\tabcolsep}{4pt}
\pgfplotstabletypeset[
  columns={label,nh,beerh,pptransh,deltah,n,beer,pptrans,delta},
  columns/label/.style={string type, column name={construct}, column type={l}},
  columns/nh/.style={column name={goals}, int detect,
                     /pgf/number format/1000 sep={\,}},
  columns/beerh/.style={column name={\beer{} \%}, fixed, fixed zerofill,
                        precision=1},
  columns/pptransh/.style={column name={ppTrans \%}, fixed, fixed zerofill,
                           precision=1},
  columns/deltah/.style={column name={$\Delta$}, fixed, fixed zerofill,
                         precision=1},
  columns/n/.style={column name={goals}, int detect,
                    /pgf/number format/1000 sep={\,}},
  columns/beer/.style={column name={\beer{} \%}, fixed, fixed zerofill,
                       precision=1},
  columns/pptrans/.style={column name={ppTrans \%}, fixed, fixed zerofill,
                          precision=1},
  columns/delta/.style={column name={$\Delta$}, fixed, fixed zerofill,
                        precision=1},
  every head row/.style={
    before row={\toprule & \multicolumn{4}{c}{verdict-only}
                & \multicolumn{4}{c}{strict} \\
                \cmidrule(lr){2-5}\cmidrule(l){6-9}},
    after row=\midrule},
  every last row/.style={after row=\bottomrule},
]{\byfeat}
\caption{Success rates by construct class under both scorings, in the order of
  \Cref{fig:eval-construct}.}
\label{tab:eval-construct}
\end{table}

\begin{table}[t]
\centering\small
\setlength{\tabcolsep}{4pt}
\pgfplotstabletypeset[
  columns={label,n,beer,pptrans,delta,noscript,pptkilled,beerh,pptransh},
  columns/label/.style={string type, column name={order}, column type={l}},
  columns/n/.style={column name={goals}, int detect,
                    /pgf/number format/1000 sep={\,}},
  columns/beer/.style={column name={\beer{} \%}, fixed, fixed zerofill,
                       precision=1},
  columns/pptrans/.style={column name={ppTrans \%}, fixed, fixed zerofill,
                          precision=1},
  columns/delta/.style={column name={$\Delta$}, fixed, fixed zerofill,
                        precision=1},
  columns/noscript/.style={column name={no script \%}, fixed, fixed zerofill,
                           precision=1},
  columns/pptkilled/.style={column name={ppT killed \%}, fixed, fixed zerofill,
                            precision=1},
  columns/beerh/.style={column name={\beer{} \%}, fixed, fixed zerofill,
                        precision=1},
  columns/pptransh/.style={column name={ppTrans \%}, fixed, fixed zerofill,
                           precision=1},
  every head row/.style={
    before row={\toprule & & \multicolumn{5}{c}{strict}
                & \multicolumn{2}{c}{verdict-only} \\
                \cmidrule(lr){3-7}\cmidrule(l){8-9}},
    after row=\midrule},
  every last row/.style={after row=\bottomrule},
]{\byorder}
\caption{Success rates by type order under both scorings.}
\label{tab:eval-order}
\end{table}

